%%%%%%%%%%%%%%%%%%%%%%%%%
\spellentry{Teleportation Circle}
%%%%%%%%%%%%%%%%%%%%%%%%%

\tagline{Conjuration (Teleportation)}
\begin{description*}
\item[Level:] Sor/Wiz 9
\item[Casting Time:] 10 minutes
\item[Range:] 0 ft.
\item[Effect:] 5-ft.-radius circle that teleports those who activate it
\item[Duration:] 10 min./level (D)
\item[Saving Throw:] None
\item[Spell Resistance:] Yes
\item[Components:] V, M
\item[Arcane Material Component:] Amber dust to cover the area of the circle (cost 1,000
\end{description*}

You create a circle on the floor or other horizontal surface that teleports, as
\textit{greater teleport}, any creature who stands on it to a designated spot.
Once you designate the destination for the circle, you can't change it. The spell
fails if you attempt to set the circle to teleport creatures into a solid object,
to a place with which you are not familiar and have no clear description, or to
another plane.
The circle itself is subtle and nearly impossible to notice. If you intend to keep
creatures from activating it accidentally, you need to mark the circle in some
way.
\textit{Teleportation circle} can be made permanent with a \textit{permanency} spell.
A permanent \textit{teleportation circle} that is disabled becomes inactive for
10 minutes, then can be triggered again as normal.
Note: Magic traps such as \textit{teleportation circle} are hard to detect
and disable. A rogue (only) can use the Search skill to find the circle and Disable
Device to thwart it. The DC in each case is 25 + spell level, or 34 in the case
of \textit{teleportation circle}.
gp).